\wheader{}{10}{4}{4}{100pt}

Pendahuluan Jilid 2 \vtmpb{3.3.03}\pagereference{10}{}

%wheader parameters:  #1 new paragraph name
 % #2#3#4  \vskip parameters if page break not forced (no. of points)
 % #5  spare height demanded if page break not to be forced
\wheader{}{10}{4}{4}{100pt}

 *Bab 21: Taksonomi ruang ukur

\chapintrosection{17.1.15}{12}{}

\section{211}{Definisi}{20.11.03}{12}{}
{Ruang ukur lengkap, berhingga total, $\sigma$-hingga, dapat dilokalkan
secara ketat, semihingga, dapat dilokalkan, dan ditentukan secara lokal;
atom; hubungan elementer; ukuran terhitung-koterhitung.}

\section{212}{Ruang lengkap}{10.9.04}{17}{}
{Fungsi terukur dan terintegralkan pada ruang lengkap;
pelengkapan suatu ukuran.}

\section{213}{Ruang semihingga, ditentukan secara lokal, dan dapat
dilokalkan}{13.9.13}{23}{}
{Integrasi pada ruang semihingga; versi c.l.d.;\ selubung
terukur; karakterisasi sifat dapat dilokalkan dan dapat dilokalkan secara
ketat.}

\section{214}{Subruang}{22.5.09}{34}{}
{Ukuran subruang pada himpunan bagian sebarang; integrasi; jumlah
langsung ruang ukur; *memperluas ukuran ke keluarga himpunan yang
terurut baik.}

\section{215}{Ruang $\sigma$-hingga dan prinsip penghabisan}
{13.11.13}{44}{}
{Prinsip penghabisan; karakterisasi sifat $\sigma$-hingga;
teorema nilai antara untuk ukuran tanpa atom.}

\section{*216}{Contoh-contoh}{25.9.04}{48}{}
{Ruang lengkap yang dapat dilokalkan tetapi tidak ditentukan secara lokal;
ruang lengkap yang ditentukan secara lokal tetapi tidak dapat dilokalkan;
ruang lengkap yang ditentukan secara lokal dan dapat dilokalkan, tetapi
tidak dapat dilokalkan secara ketat.}

\wheader{}{10}{4}{4}{100pt}

Bab 22:  Teorema Dasar Kalkulus

\chapintrosection{24.2.14}{55}{}

\section{221}{Teorema Vitali di $\Bbb{R}$}{2.6.03}{55}{}
{Teorema Vitali untuk selang di $\Bbb{R}$.}

\section{222}{Mendiferensialkan integral tak tentu}{20.11.03/18.10.04}
{58}{}
{Fungsi monoton dapat didiferensialkan hampir di mana-mana, dan
turunannya terintegralkan; $\bover{d}{dx}\int_a^xf=f$ hampir di mana-mana;
*Teorema Denjoy-Young-Saks.}

\section{223}{Teorema-teorema kerapatan Lebesgue}{9.9.04}{66}{}
{$f(x)=\lim_{h{\downarrow}0}\bover1{2h}\int_{x-h}^{x+h}f$ hampir di
mana-mana\ ($x$); titik kerapatan;
$\lim_{h{\downarrow}0}\bover1{2h}\int_{x-h}^{x+h}|f-f(x)|=0$ hampir di
mana-mana\ ($x$); himpunan Lebesgue suatu fungsi.}

\section{224}{Fungsi bervariasi terbatas}{29.9.04}{70}{}
{Variasi suatu fungsi; selisih fungsi-fungsi monoton;
jumlah dan hasil kali, limit, kekontinuan, dan keterdiferensialan bagi
fungsi bervariasi\
terbatas; suatu ketaksamaan bagi ${\int}f{\times}g$.}

\section{225}{Fungsi kontinu mutlak}{16.8.15}{79}{}
{Kekontinuan mutlak integral tak tentu; fungsi kontinu mutlak
pada $\Bbb{R}$; integrasi parsial; fungsi semikontinu bawah;
*citra langsung himpunan terabaikan; fungsi Cantor.}

\section{*226}{Dekomposisi Lebesgue suatu fungsi bervariasi
terbatas}{6.11.13}{89}{}
{Jumlah pada himpunan indeks sebarang; fungsi saltus; dekomposisi
Lebesgue.}

\wheader{}{10}{4}{4}{100pt}

 Bab 23: Teorema Radon-Nikod\'ym

\chapintrosection{17.11.04}{96}{}

\section{231}{Fungsional aditif terhitung}{25.8.15}{96}{}
{Fungsional aditif dan aditif terhitung; dekomposisi Jordan dan Hahn.}

\section{232}{Teorema Radon-Nikod\'ym}{5.6.02}{100}{}
{Fungsional aditif kontinu mutlak dan kontinu sejati; fungsional
kontinu sejati merupakan integral tak tentu; *dekomposisi Lebesgue
suatu fungsional aditif terhitung.}

\section{233}{Ekspektasi bersyarat}{16.6.02}{109}{}
{Subaljabar-$\sigma$; ekspektasi bersyarat fungsi terintegralkan;
fungsi konveks; ketaksamaan Jensen.}

\section{234}{Operasi pada ukuran}{11.4.09}{117}{}
{Fungsi pelestari ukuran invers; ukuran citra; jumlah ukuran;
ukuran integral tak tentu; pengurutan ukuran.}

\section{235}{Transformasi terukur}{30.3.03}{127}{}
{Rumus ${\int}g(y)\nu(dy)={\int}J(x)g(\phi(x))\mu(dx)$;
syarat penerapan secara terperinci; fungsi \imp;\ malapetaka ukuran
citra; penggunaan Teorema Radon-Nikod\'ym.}

\wheader{}{10}{4}{4}{100pt}

 Bab 24:  Ruang fungsi

\chapintrosection{15.11.13}{138}{}  

\section{241}{$\eusm{L}^0$ dan $L^0$}{6.11.03}{138}{}
{Struktur linear, tertib, dan multiplikatif $L^0$; kelengkapan Dedekind
dan sifat dapat dilokalkan; aksi fungsi Borel.}

\section{242}{$L^1$}{19.11.03}{146}{}
{Kekisi bernorma $L^1$; integrasi sebagai fungsional linear;
kelengkapan dan kelengkapan Dedekind; Teorema Radon-Nikod\'ym dan
ekspektasi bersyarat; fungsi konveks; subruang padat.}

\section{243}{$L^{\infty}$}{30.4.04}{156}{}{Kekisi bernorma
$L^{\infty}$; kelengkapan norma; dualitas
antara $L^1$ dan $L^{\infty}$; sifat dapat dilokalkan, kelengkapan
Dedekind, dan identifikasi $L^{\infty}\cong(L^1)^*$.}

\section{244}{$L^p$}{6.3.09}{164}{}
{Kekisi bernorma $L^p$, untuk $1<p<\infty$; ketaksamaan H\"older;
kelengkapan dan kelengkapan Dedekind;
$(L^p)^*\cong{L^q}$; ekspektasi bersyarat; *kekonveksan seragam.}

\section{245}{Konvergensi dalam ukuran}{25.3.06}{179}{}
{Topologi konvergensi (lokal) dalam ukuran pada $L^0$;
konvergensi titik demi titik; sifat dapat dilokalkan dan kelengkapan
Dedekind; penanaman $L^p$ ke dalam $L^0$; konvergensi
$\|\,\|_1$-\vthsp{} dan konvergensi dalam
ukuran; ruang $\sigma$-hingga, kemetrikan, dan konvergensi
sekuensial.}

\section{246}{Keterintegralan seragam}{17.11.06}{190}{}
{Himpunan terintegralkan secara seragam dalam $\eusm{L}^1$ dan $L^1$;
sifat elementer; karakterisasi melalui barisan saling lepas; $\|\,\|_1$
dan konvergensi dalam ukuran pada himpunan terintegralkan secara seragam.}

\section{247}{Kekompakan lemah dalam $L^1$}{26.8.13}{198}{}
{Suatu himpunan bagian $L^1$ terintegralkan secara seragam jika dan hanya
jika relatif kompak lemah.}

\wheader{}{10}{4}{4}{100pt}

 Bab 25:  Ukuran produk

\chapintrosection{31.5.03}{204}{}

\section{251}{Produk berhingga}{10.11.06}{204}{}
{Produk primitif dan c.l.d.;\ sifat-sifat dasar; ukuran Lebesgue
pada $\Bbb{R}^{r+s}$ sebagai ukuran produk; produk jumlah langsung dan
subruang;\ versi c.l.d.}

\section{252}{Teorema Fubini}{6.12.07}{212}{}
{Kapan ${\iint}f(x,y)dxdy$ dan ${\int}f(x,y)d(x,y)$ sama; ukuran
himpunan ordinat; *volume bola di $\Bbb{R}^r$.}

\section{253}{Produk tensor}{18.4.08}{237}{}
{Operator bilinear; operator bilinear $L^1(\mu){\times}L^1(\nu){\to}W$
dan operator linear $L^1(\mu\times\nu){\to}W$; operator bilinear
positif dan pengurutan pada $L^1(\mu\times\nu)$;
ekspektasi bersyarat; integral atas.}

\section{254}{Produk tak hingga}{23.2.16}{248}{}
{Produk keluarga sebarang ruang probabilitas; sifat-sifat
dasar; fungsi \imp;\ ukuran biasa pada $\{0,1\}^I$;
$\{0,1\}^{\Bbb{N}}$ isomorfik, sebagai ruang ukur, dengan $[0,1]$; subruang
berukuran luar penuh; himpunan yang ditentukan oleh koordinat dalam
suatu himpunan bagian dari himpunan indeks; hukum asosiatif tergeneralisasi
untuk produk ukuran; subproduk sebagai ukuran citra; pemfaktoran fungsi
melalui subproduk; ekspektasi bersyarat pada subaljabar yang bersesuaian
dengan subproduk; produk ruang yang dapat dilokalkan; produk ruang tanpa
atom.}

\section{255}{Konvolusi fungsi}{3.7.08}{266}{}
{Geseran di $\Bbb{R}^2$ sebagai automorfisme ruang ukur;
konvolusi fungsi pada $\Bbb{R}$;
${\int}h\times(f*g)={\int}h(x+y)f(x)g(y)d(x,y)$; $f*(g*h)=(f*g)*h$;
$\|f*g\|_1\le
\|f\|_1\|g\|_1$;  %this break required by  mtchrefs.for
grup $\Bbb{R}^r$ dan $\ocint{-\pi,\pi}$.}

\section{256}{Ukuran Radon pada $\Bbb{R}^r$}{6.8.15}{277}{}
{Definisi ukuran Radon pada $\Bbb{R}^r$; pelengkapan ukuran Borel;
keterukuran Lusin; ukuran citra;
produk dua ukuran Radon; fungsi semikontinu.}

\section{257}{Konvolusi ukuran}{14.8.13}{285}{}
{Konvolusi ukuran Radon berhingga total pada $\Bbb{R}^r$;
${\int}h\,d(\nu_1*\nu_2)={\iint}h(x+y)\nu_1(dx)\nu_2(dy)$;
$\nu_1*(\nu_2*\nu_3)=(\nu_1*\nu_2)*\nu_3$; konvolusi dan
turunan Radon-Nikod\'ym.}

\wheader{}{10}{4}{4}{100pt}

 Bab 26:  Perubahan variabel dalam integral

\chapintrosection{5.9.03}{288}{}

\section{261}{Teorema Vitali di $\BbbR^r$}{11.12.12}{288}{}
{Teorema Vitali untuk bola di $\BbbR^r$; Teorema Kerapatan
Lebesgue; himpunan Lebesgue.}

\section{262}{Fungsi Lipschitz dan terdiferensialkan}{11.8.15}{296}{}
{Fungsi Lipschitz; sifat-sifat elementer; fungsi terdiferensialkan
dari $\BbbR^r$ ke $\BbbR^s$; keterdiferensialan dan turunan parsial;
menghampiri fungsi terdiferensialkan dengan fungsi afin sesepenggal;
*Teorema Rademacher.}

\section{263}{Transformasi terdiferensialkan di $\BbbR^r$}{4.4.13}{308}{}
{Dalam rumus ${\int}g(y)dy={\int}J(x)g(\phi(x))dx$, tentukan $J$ ketika
$\phi$ (i) linear (ii) terdiferensialkan; syarat penerapan secara
terperinci; koordinat polar; kasus $\phi$ tidak injektif;
kasus satu dimensi.}

\section{264}{Ukuran Hausdorff}{12.5.03}{320}{}
{Ukuran Hausdorff berdimensi $r$ pada $\BbbR^s$; himpunan Borel
bersifat terukur; fungsi Lipschitz; jika $s=r$, diperoleh suatu kelipatan
ukuran Lebesgue; *ukuran Cantor sebagai ukuran Hausdorff.}

\section{265}{Ukuran permukaan}{3.9.13}{330}{}
{Ukuran Hausdorff ternormalisasi; aksi operator linear dan
fungsi terdiferensialkan; ukuran permukaan pada bola.}

\section{*266}{Ketaksamaan Brunn-Minkowski}{28.1.09}{338}{}
{Rataan aritmetika dan geometri; tutupan esensial; ketaksamaan
Brunn-Minkowski.}

\wheader{}{10}{4}{4}{100pt}

 Bab 27:  Teori probabilitas

\chapintrosection{26.8.13}{343}{}

\section{271}{Distribusi}{11.12.08}{344}{}
{Terminologi; distribusi sebagai ukuran Radon; fungsi distribusi;
kerapatan; transformasi peubah acak; *fungsi distribusi dan konvergensi
dalam ukuran.}

\section{272}{Kebebasan}{3.4.09}{351}{}
{Keluarga peubah acak yang saling bebas; karakterisasi
kebebasan; distribusi gabungan keluarga bebas (berhingga), dan
ukuran produk; hukum nol-satu; $\Expn(X{\times}Y)$, $\Var(X+Y)$;
distribusi suatu jumlah sebagai konvolusi distribusi; ketaksamaan
Etemadi; *ketaksamaan Hoeffding.}

\section{273}{Hukum kuat bilangan besar}{2.12.09}{364}{}
{$\bover1{n+1}\sum_{i=0}^nX_i{\to}0$ hampir di mana-mana\ jika $X_n$
saling bebas dengan ekspektasi nol dan salah satu dari (i)
$\sum_{n=0}^{\infty}\bover1{(n+1)^2}$
\discretionary{}{}{}$\Var(X_n)<\infty$ atau (ii)
$\sum_{n=0}^{\infty}\Expn(|X_n|^{1+\delta})<\infty$ untuk suatu $\delta>0$
atau (iii) $X_n$ berdistribusi identik.}

\section{274}{Teorema Limit Pusat}{13.4.10}{376}{}
{Peubah acak berdistribusi normal; syarat Lindeberg untuk
Teorema Limit Pusat; akibat-akibat; menaksir
$\int_{\alpha}^{\infty}e^{-x^2/2}dx$.}

\section{275}{Martingal}{3.12.12}{388}{}
{Barisan aljabar-$\sigma$, dan martingal yang teradaptasi padanya;
persilangan naik; Teorema Konvergensi Martingal Doob; keterintegralan
seragam, konvergensi $\|\,\|_1$-\vthsp{} dan martingal sebagai
barisan ekspektasi bersyarat; martingal balik; waktu henti.}

\section{276}{Barisan selisih martingal}{16.4.13}{399}{}
{Barisan selisih martingal; hukum kuat
bilangan besar untuk barisan selisih martingal; Teorema Koml\'os.}

\wheader{}{10}{4}{4}{100pt}

 Bab 28:  Analisis Fourier

\chapintrosection{17.1.15}{408}{}

\section{281}{Teorema Stone-Weierstrass}{4.12.12}{408}{}
{Menghampiri suatu fungsi pada himpunan kompak dengan anggota suatu
kekisi atau aljabar fungsi yang diberikan; kasus riil dan kompleks; hampiran
dengan polinomial dan fungsi trigonometri; Teorema Ekuistribusi Weyl
di $[0,1]^r$.}

\section{282}{Deret Fourier}{24.9.09}{419}{}
{Jumlah Fourier dan Fej\'er; kernel Dirichlet dan Fej\'er;
Lema Riemann-Lebesgue; konvergensi seragam jumlah Fej\'er dari
fungsi kontinu; konvergensi hampir di mana-mana\ dan $\|\,\|_1$ dari jumlah
Fej\'er suatu fungsi
terintegralkan; konvergensi $\|\,\|_2$ dari jumlah Fourier suatu fungsi
terintegralkan-kuadrat; konvergensi jumlah Fourier suatu fungsi
terdiferensialkan atau bervariasi\
terbatas;
konvolusi dan koefisien Fourier.}

\section{283}{Transformasi Fourier I}{31.3.13}{438}{}
{Transformasi Fourier dan transformasi Fourier invers; sifat-sifat
elementer; $\int_0^{\infty}\bover1x\sin{x}\,dx
\ifdim\pagewidth>467pt\penalty-100\fi
=\bover12\pi$; rumus
$\varhatf\varcheck{\phantom{h}}=f$ untuk $f$ yang terdiferensialkan dan
bervariasi terbatas;\ konvolusi; $e^{-x^2/2}$;
${\int}f\times\varhat{g}=\int\varhatf{\times}g$.}

\section{284}{Transformasi Fourier II}{30.8.13}{453}{}
{Fungsi uji; $\varhat{h}\varcheck{\phantom{h}}=h$; fungsi
tempered; fungsi tempered yang merepresentasikan transformasi satu sama lain;
konvolusi; fungsi terintegralkan-kuadrat; fungsi delta Dirac.}

\section{285}{Fungsi karakteristik}{18.9.14}{470}{}
{Fungsi karakteristik suatu distribusi; peubah acak yang saling
bebas; distribusi normal; topologi samar pada ruang
distribusi, dan konvergensi sekuensial fungsi karakteristik;
Teorema Poisson; konvolusi distribusi.}

\section{*286}{Teorema Carleson}{30.3.16}{485}{}
{Teorema Maksimal Hardy-Littlewood; bukti Lacey-Thiele untuk
Teorema Carleson bagi fungsi terintegralkan-kuadrat pada $\Bbb{R}$ dan
$\ocint{-\pi,\pi}$.}

\wheader{}{10}{4}{4}{100pt}

 Lampiran Jilid 2

\chapintrosection{3.8.15}{518}{}

\section{2A1}{Teori himpunan}{20.1.13}{518}{}
{Himpunan terurut; rekursi transfinit; ordinal; ordinal awal;
Teorema Schr\"oder-\vthsp{}Bernstein; filter; Aksioma Pilihan;
Teorema Pengurutan Baik Zermelo; Lema Zorn; ultrafilter; suatu teorema
dalam kombinatorika.}

\section{2A2}{Topologi ruang Euklides}{30.11.09}{524}{}
{Tutupan; fungsi kontinu; himpunan kompak; himpunan terbuka di $\Bbb{R}$.}

\section{2A3}{Topologi umum}{25.7.07}{527}{}
{Topologi; fungsi kontinu; topologi subruang;
tutupan dan interior;
topologi Hausdorff; pseudometrik; konvergensi barisan;
ruang kompak; titik gugus barisan; konvergensi filter;
limit superior dan limit inferior; topologi produk; himpunan padat.}

\section{2A4}{Ruang bernorma}{4.3.14}{536}{}
{Ruang bernorma; subruang linear; ruang Banach;
operator linear berbatas; ruang dual; memperluas suatu operator linear
dari subruang padat; aljabar bernorma.}

\section{2A5}{Ruang topologis linear}{13.11.07}{539}{}
{Ruang topologis linear; topologi yang didefinisikan oleh fungsional;
himpunan konveks; kelengkapan; topologi lemah.}

\section{2A6}{Faktorisasi matriks}{10.11.14}{543}{}
{Determinan; keluarga ortonormal; $T=PDQ$ dengan $D$
diagonal serta $P$, $Q$ ortogonal.}

\wheader{}{10}{4}{4}{100pt}

Konkordansi \pagereference{545}{}

\medskip

Referensi untuk Jilid 2 \vtmpb{17.12.12}\pagereference{547}{}

\medskip

Indeks Jilid 1 dan 2

\qquad Topik dan hasil utama \pagereference{549}{}

\qquad Indeks umum \pagereference{553}{}

%570 pages
